% !PRINTED 337
\weberpage{337}
Wir setzen demnach diese Determinante
\weq{2}{\Delta_{\lambda,\mu}(x) = Lf + M\Delta,}

worin $L$ und $M$ binäre Formen 3\textsuperscript{ten} Grades von
$\lambda, \mu$ sind.

Wenn wir eine beliebige lineare Transformation auf die Variablen $x$
anwenden, so ist, wegen der Invarianten-Eigenschaft der Function $\Delta$
\weq{3}{\Delta_{\lambda,\mu}(y,b) = r^2\Delta_{\lambda,\mu\,r^2}(x,a),}

und daraus folgt, dass die Coëfficienten der einzelnen Potenzen und
Producte von $\lambda, \mu$ Covarianten der ursprünglichen Form sind, und
die Coëfficienten $L, M$ von $f$ und $\Delta$ in diesen Ausdrücken, die
rational von den Coëfficienten $a$ der Form $f$ abhängen, sind daher
\gesperrt{Invarianten} dieser Form. Wir haben darin ein Mittel, diese Invarianten
thatsächlich zu berechnen.

Bezeichnen wir nämlich mit $a_i, A_i, B_i$ die Coëfficienten der
einzelnen Potenzen und Producte der Variablen $x$ in $f, \Delta$ und
$\Delta_{\lambda,\mu}$, so ergiebt sich aus (2) ein System von zehn
Gleichungen:
\[
L a_i + M A_i = B_i,
\]
woraus $L$ und $M$ bestimmt werden können. Man erhält, wenn $i, k$ zwei
verschiedene Indices sind,
\weq{4}{\begin{aligned}
L\,(a_i A_k - a_k A_i) &= B_i A_k - B_k A_i\\
M\,(a_i A_k - a_k A_i) &= -B_i a_k + B_k a_i.
\end{aligned}}

Daraus schliesst man noch, dass $L, M$ \gesperrt{ganze} Functionen der
Coëfficienten $a_i$ sind. Denn in (4) sind die rechten Seiten ganze
Functionen, und wenn also $L$ oder $M$, in der einfachsten Gestalt
dargestellt, noch einen Nenner hätten, so müssten die sämmtlichen
Determinanten $a_iA_k - a_kA_i$, welches ganze Functionen
4\textsuperscript{ten} Grades der $a$ sind, einen gemeinschaftlichen
Theiler haben. Dass dies aber nicht möglich ist, kann man an irgend
einem speciellen Falle nachweisen, z.\,B. an der Form
\[
f = 3x_3(\alpha x_1^2 + \beta x_2^2) + a x_1^3 + b x_2^3,
\]
deren Hesse'sche Form:
\[
-\beta^2(\alpha x_3 + a x_1)x_2^2 - \alpha^2(\beta x_3 + b x_2)x_1^2
\]
ist. Hier kommen unter den $a_iA_k - a_kA_i$ z.\,B. die beiden
Functionen $a^2\beta^2, b^2\alpha^2$ vor, die keinen gemeinschaftlichen
Theiler haben.
